% Part: first-order-logic
% Chapter: models-theories
% Section: introduction

\documentclass[../../../include/open-logic-section]{subfiles}

\begin{document}

\olfileid{fol}{mat}{int}
\olsection{Pambuka}

\begin{explain}
Pangembangan metode aksiomatis iku salah siji prestasi
wigati ing sajarah ilmu, lan wigatine mligi ing sajarah
matematika. Ngembangake sawenehing bidang kanthi cara
aksiomatis mbutuhake panjlentrehan akeh pitakon: Apa
kang dadi pokok kajian bidang kasebut? Apa konsep-konsep
kang paling dhasar? Kepriye sesambungane? Apa kabeh
konsep ing bidang kasebut bisa ditemtokake lumantar
konsep-konsep dhasar iki? Apa hukum kang dituruti,
lan kang kudu dituruti, dening konsep-konsep kasebut?

Metode aksiomatis lan logika cocog siji lan sijine.
Logika formal nyedhiyakake piranti kanggo ngrumusake
teori aksiomatis, kanggo mbuktekake teorema saka aksioma
teori kanthi cara kang ditemtokake kanthi cetha,
lan kanggo nliti sipat-sipat kabeh sistem kang nyukupi
aksioma kasebut kanthi sistematis.
\end{explain}

\begin{defn}
Himpunan !!{sentence}s~$\Gamma$ diarani \emph{katutup}
yen lan mung yen, kapan wae $\Gamma \Entails !A$,
mula $!A \in \Gamma$. \emph{Panutupan} saka himpunan
!!{sentence}s~$\Gamma$ yaiku
$\Setabs{!A}{\Gamma \Entails !A}$.

Kita kandha yen~$\Gamma$ \emph{diaksiomatisasi dening}
himpunan ukara~$\Delta$ yen $\Gamma$ iku panutupan
saka~$\Delta$.
\end{defn}

\begin{explain}
Kita bisa nganggep teori aksiomatis minangka himpunan
ukara kang diaksiomatisasi dening himpunan aksiomane~$\Delta$.
Kanthi tembung liya, yen kita nduweni basa logika tataran kapisan
kang ngemot simbol nonlogis kanggo konsep-konsep dhasar
ing ilmu kang arep dikembangake kanthi aksiomatis,
bebarengan karo himpunan !!{sentence}s kang ngandharake
hukum-hukum dhasar ilmu kasebut, teori iku bisa kita
wakili kanthi kabeh !!{sentence}s ing basa kasebut
kang dadi akibat saka aksioma. Cakupane wiwit saka
tuladha prasaja kanthi mung siji konsep dhasar lan
aksioma kang prasaja, kayata teori urutan parsial,
nganti teori rumit kayata mekanika Newton.

Kasunyatan logis kang ndadekake cara formal marang
metode aksiomatis iki wigati banget yaiku ing ngisor iki.
Upamakna $\Gamma$ iku sistem aksioma kanggo sawenehing
teori, yaiku himpunan ukara.
\begin{enumerate}
\item Kita bisa nyatakake kanthi cetha kapan sistem
  aksioma nyakup kelas !!{structure}s kang dikarepake.
  Tegese, yen kita nggatekake kelas !!{structure}s tartamtu,
  sistem aksioma~$\Gamma$ kasil nyakup kelas kasebut
  yen lan mung yen !!{structure}s ing kelas iku persis
  struktur~$\Struct M$ kang ndadekake
  $\Sat{M}{\Gamma}$.
\item Kita bisa gagal ing bab iki amarga ana $\Struct M$
  kang ndadekake $\Sat{M}{\Gamma}$, nanging
  $\Struct M$ dudu salah siji !!{structure}s kang
  dikarepake. Iki bisa ndadekake kita nambah aksioma
  kang ora bener ing~$\Struct M$.
\item Yen kita paling ora kasil njaga supaya $\Gamma$
  bener ing kabeh !!{structure}s kang dikarepake,
  mula ukara~$!A$ bener ing kabeh !!{structure}s
  kang dikarepake kapan wae $\Gamma \Entails !A$.
  Dadi kita bisa nggunakake piranti logis
  (kayata metode !!{derivation}) kanggo nuduhake
  yen ukara bener ing kabeh !!{structure}s
  kang dikarepake mung kanthi nuduhake yen ukara
  kasebut dadi akibat saka aksioma.
\item Kadhang kita ora miwiti saka !!{structure}s
  kang wis dibayangake, nanging saka aksioma dhewe:
  kita miwiti saka sawetara konsep dhasar kang
  dikarepake nyukupi hukum tartamtu kang kita
  rumusake ing sistem aksioma. Salah siji bab
  kang arep langsung kita priksa yaiku apa
  aksioma-aksioma kasebut ora padha kontradiksi:
  yen kontradiksi, ora ana konsep kang bisa
  manut hukum kasebut, lan teori kang kita
  susun ora koheren. Kita bisa mriksa yen
  prakara iki ora kedadeyan kanthi nemokake
  model saka~$\Gamma$. Yen teori kita nduweni
  model, kita bisa nggunakake metode logis
  kanggo nliti model-model kasebut, lan uga
  kanggo mbangun model.
\item Kamardikan aksioma uga pitakon wigati.
  Bisa wae salah siji aksioma nyatane dadi
  akibat saka aksioma-aksioma liyane,
  mula aksioma iku keluwihan. Kita bisa
  mbuktekake yen aksioma $!A$ ing $\Gamma$
  keluwihan kanthi mbuktekake $\Gamma \setminus
  \{!A\} \Entails !A$. Kita uga bisa
  mbuktekake yen aksioma iku ora keluwihan
  kanthi nuduhake yen $(\Gamma \setminus \{!A\})
  \cup \{\lnot !A\}$ bisa dicukupi. Upamane,
  kanthi cara iki postulat paralel kabukten
  mardika saka aksioma geometri liyane.
\item Pitakon wigati liyane yaiku apa konsep
  bisa ditemtokake ing sawijining teori:
  pilihan basa nemtokake apa kang kalebu
  model-model teori. Nanging ora saben
  aspek teori kudu diwakili dhewe-dhewe
  ing modele. Upamane, saben relasi urutan
  $\le$ nemtokake relasi urutan ketat~$<$;
  yen salah siji wis ana, liyane bisa kita
  temtokake. Dadi model teori kang ngemot
  urutan kaya mangkono ora kudu \emph{uga}
  ngemot relasi urutan ketat kang cocog.
  Kapan, sacara umum, siji relasi bisa
  ditemtokake lumantar relasi liyane?
  Kapan siji relasi ora bisa ditemtokake
  lumantar relasi liyane (lan mula kudu
  ditambahake dadi konsep dhasar ing basa)?
\end{enumerate}
\end{explain}


\end{document}
